\section{Data}
\label{sec:data}

The system relies on three distinct data assets: (i) raw course material that backs the
retrieval index, (ii) public machine-learning question--answer (QA) collections used as
auxiliary supervision, and (iii) a synthetic task-specific corpus generated from the course
material to teach the model the four output formats. The synthetic corpus is described in
Section~\ref{sec:training_data} together with the model implementation, since its construction
is tightly coupled to the training objective. This section describes the first two and the
indexing pipeline that turns the course material into a searchable RAG index.

\subsection{Course material}

The primary source is \textit{Deep Learning with Python} (3rd ed.) by F.\ Chollet
\cite{chollet2024dlwithpython}, the textbook used in DAT255. The chapters were converted to
Markdown so that the natural document hierarchy --- chapter, section, subsection --- is
preserved as explicit \texttt{\#}, \texttt{\#\#}, \texttt{\#\#\#} headers. Markdown was chosen
over plain PDF text extraction because the heading structure is a reliable signal of topical
boundaries. Cutting on header boundaries produces chunks that correspond to a single coherent
idea, which both improves retrieval relevance and avoids splitting an explanation halfway
through.

\subsection{Public QA datasets}

Two public QA collections from HuggingFace,
\texttt{prsdm/\allowbreak Machine-\allowbreak Learning-\allowbreak QA-\allowbreak dataset}
and
\texttt{win-wang/\allowbreak Machine\_\allowbreak Learning\_\allowbreak QA\_\allowbreak Collection},
were used as a general machine-learning corpus. Together they contain several thousand
short-form QA pairs covering supervised learning, neural networks, optimisation, and related
topics. Their value is breadth: they expose the model to a wider distribution of phrasings
than the course book alone and reduce overfitting to the textbook style. Their main
limitation is that they are not specific to the DAT255 syllabus; they are therefore used
only as an auxiliary signal and do not appear in the test split.

\subsection{RAG indexing pipeline}

The RAG pipeline is a three-stage process: chunking, embedding, and indexing
(Figure~\ref{fig:rag_pipeline}). Each stage is a separate Python script with explicit
inputs and outputs, which makes it cheap to re-run a single stage in isolation when the
upstream data changes.

\begin{figure}[h]
\centering
\begin{tikzpicture}[node distance=0.45cm and 0.7cm]

    \node[data] (md)        {\texttt{.md} files \\ \footnotesize(course book)};
    \node[block, right=of md] (chunker) {\texttt{chunker.py} \\ \footnotesize header-aware split};
    \node[data, right=of chunker] (chunks) {\texttt{chunks.json} \\ \footnotesize 1\,960 entries};
    \node[block, right=of chunks] (embed) {\texttt{embedder.py} \\ \footnotesize bge-base 768d};
    \node[data, right=of embed] (idx) {\texttt{rag\_index/} \\ \footnotesize HNSW + meta};

    \node[block, below=1cm of embed] (retr) {\texttt{retriever.py}};
    \node[data, left=1.6cm of retr] (q) {Query};
    \node[data, right=1.4cm of retr] (topk) {Top-$k$ chunks};

    \draw[arrow] (md) -- (chunker);
    \draw[arrow] (chunker) -- (chunks);
    \draw[arrow] (chunks) -- (embed);
    \draw[arrow] (embed) -- (idx);
    \draw[arrow] (idx) -- (retr);
    \draw[arrow] (q) -- (retr);
    \draw[arrow] (retr) -- (topk);

\end{tikzpicture}
\caption{RAG indexing and retrieval pipeline. The first three stages run offline; only the
retriever runs at inference time.}
\label{fig:rag_pipeline}
\end{figure}

\paragraph{Chunking.}
\texttt{chunker.py} reads a manifest of Markdown files and splits them in two passes. The
first pass walks the document tree and emits one chunk per heading section at any of four
levels (\texttt{\#}, \texttt{\#\#}, \texttt{\#\#\#}, \texttt{\#\#\#\#}). Every header section
is then re-fed through a recursive character splitter with a 600-character target size and
80-character overlap between consecutive sub-chunks; the overlap preserves context across
the boundary and reduces edge-effect retrieval misses. Sub-chunks shorter than 80 characters
(typically empty header stubs) are discarded. Each chunk is stored together with breadcrumb
metadata (\texttt{h1 > h2 > h3}) so retrieved passages can be displayed with full provenance.
The resulting corpus contains 1\,960 chunks.

\paragraph{Embedding and indexing.}
\texttt{embedder.py} encodes each chunk with the \texttt{BAAI/bge-base-en-v1.5} sentence
encoder \cite{xiao2023cpack}, a 768-dimensional model that ranks competitively on MTEB while
remaining small enough to run on a single consumer GPU. A lighter \texttt{all-MiniLM-L6-v2}
sentence-transformers model in the spirit of Sentence-BERT~\cite{reimers2019sentence} is
supported as a CPU fallback. The embeddings are stored in
an HNSW approximate-nearest-neighbour index \cite{malkov2020hnsw} via the \texttt{hnswlib}
library. HNSW was chosen over a flat index because $O(\log N)$ search keeps inference latency
flat as the corpus grows; flat search would have been adequate at the current size but would
not scale to a multi-course expansion.

\paragraph{A note on tokeniser separation.}
The retriever uses the BGE encoder's BERT-style WordPiece tokenizer internally to produce
semantic embeddings. The generator uses a separate tiktoken-based GPT-2 byte-pair encoding
(BPE) tokeniser (Section~\ref{sec:tokenizer}). The two tokenisers serve different purposes
--- the first exists to compute a similarity vector for search, the second to feed integer
IDs into the generator's embedding layer --- and are deliberately not shared.

\paragraph{Format of an indexed chunk.}
Each entry in \texttt{chunks.json} stores both the raw text and the breadcrumb that
locates it in the source book:

\begin{quote}\small\ttfamily
\{ \\
\hspace*{1em}"text": "Artificial intelligence was b...", \\
\hspace*{1em}"source": ".../ch01\_what\_is\_deep\_learning.md", \\
\hspace*{1em}"metadata": \{ \\
\hspace*{2em}"h1": "Chapter 1: What is deep learning?", \\
\hspace*{2em}"h2": "Artificial intelligence", \\
\hspace*{2em}"breadcrumb": "Deep Learning with Python > Chapter 1 > Artificial intelligence" \\
\hspace*{1em}\} \\
\}
\end{quote}
